\clearpage
\subunitheader{4.2: Do the scaling rules predict longer training?}

Do the scaling ideas from Problem 4.1 help over longer training? Use the
default model and training settings at 153.6M tokens.

The course baseline uses the default architecture and initialization.
For $\mu$P, change only the four settings in the table below, using reference
width $n_0=512$. Let $n$ be hidden width, $m=n/n_0$, and $\eta$ the peak base LR.
Let $G$ be a shared matrix of standard Gaussian draws truncated at $\pm3$,
used for both embedding and readout initialization and shared across prescriptions.
\begin{center}
\begin{tabular}{lcc}
\toprule
Setting & Course baseline (default) & $\mu$P \\
\midrule
Embedding initialization & $G/n$ & $G/n_0$ \\
Readout initialization & $G/\sqrt n$ & $G/\sqrt{n_0}$ \\
Readout forward multiplier & $1$ & $1/m$ \\
Hidden-matrix peak LR & $\eta$ & $\eta/m$ \\
\bottomrule
\end{tabular}
\end{center}
Keep all other training settings fixed, including RMSNorm $\epsilon=10^{-5}$
and gradient clipping at 1. Normalization and clipping can affect widths
differently: the baseline embedding scale $1/n$ may be comparable to
$\sqrt\epsilon$. Use feature logs and \texttt{gradients.jsonl} to report
pre-clipping gradient norms, the embedding's fraction of squared gradient
norm, and the clipping coefficient. Account for these effects when interpreting
transfer. Reuse the supplied width-512 diagnostics; optional normalization or
clipping ablations must fit within your compute budget.


\begin{problem}{From five updates to a longer training budget (recommended: 50 new runs)}
\begin{parts}
\item \textbf{Transfer the reference LR (20 initial runs).}
Predict whether Problem 1's best sampled LR at 153.6M tokens transfers to
widths 128 and 256 under both prescriptions. Choose an LR sweep for each
configuration, including the transferred source LR.
Hold data and other settings fixed across new runs; compare with Problem 1's
supplied width-512 curve.

At the source LR, plot $\alpha_{\rm upd}^{(t)}$ throughout training for both
prescriptions and all three widths. Use fixed inputs and each matrix's fan-in
for every block's attention and feedforward matrices and the readout.
Also plot $\omega_{\rm move}^{(t)}$ from Problem 4.1. How do update alignment
and readout--feature-change alignment evolve beyond the first five updates?

\item \textbf{Compare transfer and tuned performance (no new runs).}
Fit each loss--LR curve. Plot fitted optimal base LR and fitted minimum validation
loss against width for both prescriptions. Also plot the measured loss at the
transferred LR. Does the prescription with better LR transfer also achieve lower
loss after tuning? How does the comparison differ from the five-step test?

\item \textbf{Predict a held-out width (about 12 new runs).}
Hold out width 1,024. For each prescription, fit a power law to the optimal base
LRs at widths 128, 256, and 512 from part (b), and predict its optimal LR at 1,024.
Train at the predicted LR and at the directly transferred source LR. Then run a
local LR sweep at 1,024. Compare both predicted LRs with the fitted local optimum and
report their measured loss gaps relative to the best sampled target LR. Does fitting the width dependence improve on
direct transfer?

\item \textbf{Transfer across depth (18 initial runs).}
Fix width 512 and compare depths 4, 8, and 16. Test standard $\mu$P,
Depth-$\mu$P, and CompleteP using the rules from Problem 4.1, now with reference
depth eight: for layer count $L$, use relative depth $r=L/8$.
At depths 4 and 16, compare the transferred source LR with an LR sweep of
your choice. Plot fitted optimal base LR and fitted minimum
validation loss against depth, alongside the measured losses at the transferred
LR. Does a depth correction improve transfer, tuned performance, or both?

\item \textbf{Identify the training regime (diagnostics from the same runs).}
At the directly transferred LR and the best sampled LR for each configuration,
track logit RMS and normalized hidden-feature movement on fixed inputs during
the first five updates and at the end of training, as in Problem 4.1. Include both the held-out width and the depth comparison.
Include $\omega_{\rm move}^{(t)}$ at the same
checkpoints. Which alignment assumptions are supported at different widths,
depths, and stages of training?
Can a configuration with larger early transients still achieve a better final loss?
Does a shift in optimal LR alone show that the model is outside the stable
$\mu$P regime? Use your diagnostics to explain what the experiments establish.
Use the LR--WD coupling from Problem 2 to discuss how holding WD fixed might
affect the comparison.
\item \textbf{Muon and $\mu$P (optional).}
\href{https://kellerjordan.github.io/posts/muon/}{Muon} is a momentum-based
optimizer that approximately orthogonalizes the update direction for each
hidden weight matrix before applying it. Explore how Muon interacts with
$\mu$P, following the same approach as in parts (a)--(c).
This exploration is not part of the required workload; attempt it only if you
have GPU budget remaining after the required experiments.

\end{parts}
\end{problem}
